% SEGMENT OLP-0362-S01
% Part: lambda-calculus
% Chapter: syntax
% Section: alpha

\documentclass[../../../include/open-logic-section]{subfiles}

\begin{document}

\olfileid{lam}{syn}{alp}
\olsection{$\alpha$-மாற்றம்}

$\lambd[x][x]$ க்கும் $\lambd[y][y]$ க்கும் இடையிலான தொடர்பு என்ன? இரண்டும்
சமனிச் சார்பைக் குறிக்கின்றன. அவை தொடரமைப்பின்படி வெவ்வேறு சொற்களே. அவற்றின்
கட்டுண்ட மாறியின் பெயரில் மட்டுமே அவை வேறுபடுகின்றன; ஒன்றின் கட்டுண்ட மாறியை
``மறுபெயரிடுவதால்'' மற்றொன்று கிடைக்கிறது. இது \emph{$\alpha$-மாற்றம்}
எனப்படுகிறது.

\begin{defn}[கட்டுண்ட மாறியின் மாற்றம், $\aconvone$]
  சொல் $M$ ஆனது $\lambd[x][N]$ இன் ஒரு நிகழ்வைக் கொண்டிருந்து,
  $y \notin \FV{N}$ ஆகவும் $\Subst{N}{y}{x}$ வரையறுக்கப்பட்டதாகவும்
  இருந்தால், இந்த நிகழ்வை
  \begin{equation*}
    \lambd[y][\Subst{N}{y}{x}]
  \end{equation*}
  ஆல் மாற்றி $M'$ ஐப் பெறுவது \emph{கட்டுண்ட மாறியின் மாற்றம்} எனப்பட்டு,
  $M \redone[\alpha] M'$ என்று எழுதப்படுகிறது.
\end{defn}

\begin{defn}[ஒரு தொடர்பின் இசைவு]
  சொற்களின்மீதான தொடர்பு $R$ பின்வரும் நிபந்தனைகளை நிறைவுசெய்தால் அது
  \emph{இசைவானது} எனப்படுகிறது:
  \begin{enumerate}
  \item $R N N'$ எனில் $R \lambd[x][N] \lambd[x][N']$
  \item $R P P'$ எனில் $R (PQ) (P'Q)$
  \item $R Q Q'$ எனில் $R (PQ) (PQ')$
  \end{enumerate}
\end{defn}

ஆகவே வரையறையை மீண்டும் கூறுவோம்:
\begin{defn}[கட்டுண்ட மாறியின் மாற்றம், $\aconvone$]
  \emph{கட்டுண்ட மாறியின் மாற்றம்} ($\redone[\alpha]$) என்பது சொற்களின்மீதான,
  பின்வரும் நிபந்தனையை நிறைவுசெய்யும் மிகச்சிறிய இசைவான தொடர்பாகும்:
  \begin{align*}
    &\lambd[x][N] \redone[\alpha] \lambd[y][\Subst{N}{y}{x}] && \text{
      $x \neq y$, $y \notin \FV{N}$ எனில்} \\
    & &&\text{மேலும் $\Subst{N}{y}{x}$ வரையறுக்கப்பட்டிருந்தால்}
  \end{align*}
\end{defn}

இங்கு ``மிகச்சிறிய'' என்பது, இசைவாலும் கூடுதல் நிபந்தனையாலும் தேவைப்படும்
ஜோடிகளை மட்டுமே இத்தொடர்பு கொண்டிருக்கும்; வேறெதையும் கொண்டிருக்காது என்பதாகும்.
ஆகவே இத்தொடர்பைப் பின்வருமாறும் வரையறுக்கலாம்:

\begin{defn}[கட்டுண்ட மாறியின் மாற்றம், $\aconvone$] \ollabel{defn:aconvone}
  \emph{கட்டுண்ட மாறியின் மாற்றம்} ($\aconvone$) பின்வருமாறு தொகுத்தறிதல்
  முறையில் வரையறுக்கப்படுகிறது:
  \begin{enumerate}
  \item $N \aconvone N'$ எனில் $\lambd[x][N] \aconvone
    \lambd[x][N']$ \ollabel{defn:aconvone1}
  \item $P \aconvone P'$ எனில் $(PQ) \aconvone (P'Q)$ \ollabel{defn:aconvone2}
  \item $Q \aconvone Q'$ எனில் $(PQ) \aconvone (PQ')$ \ollabel{defn:aconvone3}
  \item $x \neq y$, $y \notin \FV{N}$ ஆகவும் $\Subst{N}{y}{x}$
    வரையறுக்கப்பட்டதாகவும் இருந்தால்,
    $\lambd[x][N] \redone[\alpha] \lambd[y][\Subst{N}{y}{x}]$.
    \ollabel{defn:aconvone4}
  \end{enumerate}
\end{defn}

இவ்வரையறைகள் சமானமானவை; ஆனால் நிறுவலைப் பயிற்சியாக விடுகிறோம். இனி
தொகுத்தறிதல் வரையறையைப் பயன்படுத்துவோம்.

\begin{defn}[$\alpha$-மாற்றம், $\aconv$]
  \emph{$\alpha$-மாற்றம்} ($\aconv$) என்பது $\aconvone$ ஐக் கொண்டிருக்கும்,
  சொற்களின்மீதான மிகச்சிறிய தற்சுட்டு மற்றும் கடப்புத் தொடர்பாகும்.
  % SOURCE CORRECTION TA-LAM-016: Tamil resolves the frozen prose misspellings and malformed English clauses without changing their mathematical claims.
\end{defn}

மேலேபோல, ``மிகச்சிறிய'' என்பது கடப்புத்தன்மையாலும் $\aconvone$ ஆலும்
தேவைப்படும் ஜோடிகளை மட்டுமே தொடர்பு கொண்டிருக்கும் என்பதாகும்; இது பின்வரும்
சமானமான வரையறைக்கு வழிவகுக்கிறது:

\begin{defn}[$\alpha$-மாற்றம், $\aconv$] \ollabel{defn:aconv}
  \emph{$\alpha$-மாற்றம்} ($\aconv$) பின்வருமாறு தொகுத்தறிதல் முறையில்
  வரையறுக்கப்படுகிறது:
  \begin{enumerate}
  \item $P \aconv Q$ மற்றும் $Q \aconv R$ எனில், $P \aconv R$.
    \ollabel{defn:aconv1}
  \item $P \aconvone Q$ எனில், $P \aconv Q$. \ollabel{defn:aconv2}
  \item $P \aconv P$. \ollabel{defn:aconv3}
  \end{enumerate}
\end{defn}

\begin{ex}
  $\lambd[x][f x]$ ஆனது $\lambd[y][f y]$ ஆக $\alpha$-மாறுகிறது; மறுதிசையிலும்
  அவ்வாறே. முறைசாரா விதத்தில் கூறினால், இரண்டும் ஓர் உள்ளீட்டை ஏற்று, அந்த
  உள்ளீட்டின் $f$ ஐத் திருப்பித்தரும் சார்புகள்; இங்கு அவை சூழல் மாறி~$f$ ஐச்
  சுட்டுகின்றன.

  $\lambd[x][f x]$ ஆனது $\lambd[x][g x]$ ஆக $\alpha$-மாறாது. முறைசாரா
  விதத்தில் கூறினால், அவை முறையே சூழல் மாறிகள் $f$,~$g$ ஐச் சுட்டுகின்றன;
  இதனால் அவை வெவ்வேறு சார்புகள்: $f$ உம் $g$ உம் வெவ்வேறாக உள்ள சூழல்களில்
  அவை வெவ்வேறு விதத்தில் நடந்துகொள்கின்றன.
\end{ex}

\begin{prob}
  பின்வரும் சொல் ஜோடிகள் $\alpha$-மாற்றத்தக்கவையா?
  \begin{enumerate}
  \item $\lambd[x][\lambd[y][x]]$ மற்றும் $\lambd[y][\lambd[x][y]]$
  \item $\lambd[x][\lambd[y][x]]$ மற்றும் $\lambd[c][\lambd[b][a]]$
  \item $\lambd[x][\lambd[y][x]]$ மற்றும் $\lambd[c][\lambd[b][a]]$
  \end{enumerate}
\end{prob}

\begin{lem}\ollabel{lem:fv-one}
  $P \aconvone Q$ எனில் $\FV{P} = \FV{Q}$.
\end{lem}

\begin{proof}
  $P \aconvone Q$ இன் !!{derivation} மீது தொகுத்தறிதல் முறையில் நிறுவுவோம்.
  \begin{enumerate}
  \item கடைசி விதி \olref{defn:aconvone4} எனில், $P$ ஆனது
    $\lambd[x][N]$ வடிவிலும் $Q$ ஆனது
    $\lambd[y][\Subst{N}{y}{x}]$ வடிவிலும் இருக்கும்; இங்கு $x \neq y$,
    $y \notin \FV{N}$, மேலும் $\Subst{N}{y}{x}$ வரையறுக்கப்பட்டுள்ளது.
    $x \in \FV{N}$ ஆகுமா என்பதைப் பொறுத்து நிகழ்வுகளைப் பிரிக்கிறோம்:
    % SOURCE CORRECTION TA-LAM-017: restore the FV macro omitted at five formula occurrences in the frozen source.
    \begin{enumerate}
    \item $x \in \FV{N}$ எனில்:
      \begin{align*}
        \FV{\lambd[y][\Subst{N}{y}{x}]} &=
        \FV{\Subst{N}{y}{x}} \setminus \{y\} \\
        & = ((\FV{N} \setminus \{x\}) \cup \{y\}) \setminus \{y\}
         && \text{ \olref[sub]{thm:infv} இன்படி} \\
        & = \FV{N} \setminus \{x\} \\
        & = \FV{\lambd[x][N]}
      \end{align*}
    \item $x \notin \FV{N}$ எனில்:
      \begin{align*}
        \FV{\lambd[y][\Subst{N}{y}{x}]}
        & = \FV{\Subst{N}{y}{x}} \setminus \{y\} \\
        & = \FV{N} \setminus \{x\}
         && \text{\olref[sub]{thm:notinfv} இன்படி} \\
        & = \FV{\lambd[x][N]}.
      \end{align*}
    \end{enumerate}
  \item மற்ற மூன்று நிகழ்வுகள் பயிற்சிகளாக விடப்படுகின்றன.
  \end{enumerate}
\end{proof}

\begin{prob}
  \olref[lam][syn][alp]{lem:fv-one} இன் நிறுவலை நிறைவுசெய்க.
\end{prob}

\begin{lem}\ollabel{lem:inv}
  $P \aconvone Q$ எனில் $Q \aconvone P$.
\end{lem}

\begin{proof}
  $P \aconvone Q$ இன் !!{derivation} மீது தொகுத்தறிதல் முறையில் நிறுவுவோம்.
  \begin{enumerate}
  \item கடைசி விதி \olref{defn:aconvone4} எனில், $P$ ஆனது
    $\lambd[x][N]$ வடிவிலும் $Q$ ஆனது
    $\lambd[y][\Subst{N}{y}{x}]$ வடிவிலும் இருக்கும்; இங்கு $x \neq y$,
    $y \notin \FV{N}$, மேலும் $\Subst{N}{y}{x}$ வரையறுக்கப்பட்டுள்ளது.
    முதலில், \olref[sub]{thm:clr} இன்படி $y \notin
    \FV{\Subst{N}{y}{x}}$. \olref[sub]{thm:inv} இன்படி
    $\Subst{\Subst{N}{y}{x}}{x}{y}$ வரையறுக்கப்பட்டிருப்பது மட்டுமல்லாமல்,
    அது~$N$ க்குச் சமமாகவும் உள்ளது. பின்னர் \olref{defn:aconvone4} இன்படி,
    $\lambd[y][\Subst{N}{y}{x}]
    \aconvone \lambd[x][\Subst{\Subst{N}{y}{x}}{x}{y}] =
    \lambd[x][N]$.
  \end{enumerate}
\end{proof}

\begin{prob}
  \olref[lam][syn][alp]{lem:inv} இன் நிறுவலை நிறைவுசெய்க.
\end{prob}

\begin{thm}
  $\alpha$-மாற்றம் சொற்களின்மீதான ஒரு சமானத் தொடர்பு; அதாவது, அது தற்சுட்டு,
  சமச்சீர், கடப்பு ஆகிய பண்புகளைக் கொண்டது.
\end{thm}

\begin{proof}
  \begin{enumerate}
  \item ஒவ்வொரு சொல் $M$ க்கும், கட்டுண்ட மாறிகளில் \emph{பூச்சிய} மாற்றங்கள்
    செய்து $M$ ஐ~$M$ ஆக மாற்றலாம்.
  \item கட்டுண்ட மாறிகளில் செய்யப்பட்ட தொடர் மாற்றங்களால் $P$ ஆனது~$Q$ ஆக
    $\alpha$-மாறினால், அந்த மாற்றங்களை எதிர் வரிசையில்
    (\olref{lem:inv} இன்படி) புரட்டி, $Q$ இலிருந்து~$P$ ஐப் பெறலாம்.
  \item கட்டுண்ட மாறிகளில் செய்யப்பட்ட தொடர் மாற்றங்களால் $P$ ஆனது~$Q$ ஆகவும்,
    மற்றொரு தொடர் மாற்றங்களால் $Q$ ஆனது~$R$ ஆகவும் $\alpha$-மாறினால், முதல்
    தொடரைச் செய்து பின்னர் இரண்டாம் தொடரைச் செய்வதன் மூலம் $P$ ஐ~$R$ ஆக
    மாற்றலாம்.
  \end{enumerate}
\end{proof}

இனி $M$ உம் $N$ உம் \emph{$\alpha$-சமானமானவை}, $M \aeq N$, என்று கூறுவது,
$M$ ஆனது~$N$ ஆக $\alpha$-மாறும் அப்போதும் அப்போது மட்டுமே (நாம் சற்றுமுன்
காட்டியபடி, இது $N$ ஆனது~$M$ ஆக $\alpha$-மாறும் அப்போதும் அப்போது மட்டுமே).

\begin{thm}\ollabel{thm:fv}
  $M \aeq N$ எனில், $\FV{M} = \FV{N}$.
\end{thm}

\begin{proof}
  \olref{lem:fv-one} இலிருந்து உடனடியாகக் கிடைக்கிறது.
\end{proof}

\begin{lem}\ollabel{lem:sub:R}
  $R \aeq R'$ ஆகவும் $\Subst{M}{R}{y}$ வரையறுக்கப்பட்டதாகவும் இருந்தால்,
  $\Subst{M}{R'}{y}$ வரையறுக்கப்பட்டு $\Subst{M}{R}{y}$ க்கு
  $\alpha$-சமானமாக இருக்கும்.
\end{lem}

\begin{proof}
  பயிற்சி.
\end{proof}

\begin{prob}
  \olref[lam][syn][alp]{lem:sub:R} ஐ நிறுவுக.
\end{prob}

\olref[sub]{sec} இல் சில நிகழ்வுகளில் இடைநிறுத்தல் வரையறுக்கப்படவில்லை என்பதை
நினைவுகொள்க. எனினும், சொற்கள்மீது $\alpha$-மாற்றத்தைப் பயன்படுத்தி கட்டுண்ட
மாறிகளை மறுபெயரிடுவதன் மூலம் இடைநிறுத்தலை எப்போதும் வரையறுக்கப்பட்டதாகச்
செய்யலாம். பின்வரும் தேற்றம் காட்டுவதுபோல், விளைவு $\alpha$-சமானத்தைப்
பேணுகிறது:

\begin{thm}\ollabel{thm:sub}
  எந்த $M$, $R$,~$y$ க்கும், $M \aeq M'$ ஆகவும் $\Subst{M'}{R}{y}$
  வரையறுக்கப்பட்டதாகவும் அமையும் ஓர் $M'$ உள்ளது. மேலும்,
  $M'' \aeq M$ ஆகவும் $R''$ உடன் $\Subst{M''}{R''}{y}$
  வரையறுக்கப்பட்டதாகவும், மேலும் $R'' \aeq R$ ஆகவும் அமையும் மற்றொரு ஜோடி
  இருந்தால்,
  $\Subst{M'}{R}{y} \aeq \Subst{M''}{R''}{y}$.
\end{thm}

\begin{proof}
  $M$ இன் அமைப்பின்மீது தொகுத்தறிதல் முறையில் நிறுவுவோம்:
  \begin{enumerate}
  \item $M$ ஒரு மாறி~$z$: பயிற்சி.
  \item $M$ ஆனது $\lambd[x][N]$ வடிவுடையது எனக் கொள்க. $x$, $y$
    இரண்டிலிருந்தும் வேறான, $z \notin \FV{N}$ மற்றும்
    $z \notin \FV{R}$ ஆக அமையும் ஒரு மாறி $z$ ஐத் தேர்ந்தெடுக்க.
    தொகுத்தறிதல் கருதுகோளின்படி, $N' \aeq N$ ஆகவும்
    $\Subst{N'}{z}{x}$ வரையறுக்கப்பட்டதாகவும் அமையும் ஓர் $N'$ உள்ளது.
    அப்போது \olref{defn:aconvone}\olref{defn:aconvone1} இன்படி,
    $\lambd[x][N] \aeq \lambd[x][N']$ உம் ஆகும். இப்போது
    \olref{defn:aconvone}\olref{defn:aconvone4} இன்படி,
    $\lambd[x][N'] \aeq \lambd[z][\Subst{N'}{z}{x}]$. ஏனெனில்
    $z \ne x$, $z \notin \FV{N'}$, மேலும் $\Subst{N'}{z}{x}$
    வரையறுக்கப்பட்டுள்ளது. இறுதியாக, $z \neq y$ மற்றும்
    $z \notin \FV{R}$ என்பதால்,
    $\Subst{\lambd[z][\Subst{N'}{z}{x}]}{R}{y}$ வரையறுக்கப்பட்டுள்ளது.

    இதே நிபந்தனைகளை நிறைவுசெய்யும் மற்றோர் $N''$, $R''$ இருந்தால்,
    \begin{multline*}
      \Subst{(\lambd[z][\Subst{N''}{z}{x}])}{R''}{y} =\\
    \begin{aligned}
      &= \lambd[z][\Subst{\Subst{N''}{z}{x}}{R''}{y}] \\
      &\aeq \lambd[z][\Subst{\Subst{N''}{z}{x}}{R}{y}] && \text{
        \olref{lem:sub:R} இன்படி}\\
      &\aeq\lambd[z][\Subst{\Subst{N'}{z}{x}}{R}{y}]
      && \text{தொகுத்தறிதல்
        கருதுகோளின்படி}\\
      &=\Subst{(\lambd[z][\Subst{N'}{z}{x}])}{R}{y}
    \end{aligned}
    \end{multline*}
    % SOURCE CORRECTION TA-LAM-018: the two middle comparison steps are alpha-equivalences, as the cited lemma and induction hypothesis state, rather than literal equalities.
    \item $M$ ஆனது $(PQ)$ வடிவுடையது: பயிற்சி.
  \end{enumerate}
\end{proof}

\begin{prob}
  \olref[lam][syn][alp]{thm:sub} இன் நிறுவலை நிறைவுசெய்க.
\end{prob}

\begin{cor}\ollabel{cor:sub}
  எந்த $M$, $R$,~$y$ க்கும், $M' \aeq M$, $R \aeq R'$ ஆகவும்
  $\Subst{M'}{R'}{y}$ வரையறுக்கப்பட்டதாகவும் அமையும் $M'$, $R'$ ஜோடி ஒன்று
  உள்ளது. மேலும், $M'' \aeq M$, $R'' \aeq R$ ஆகவும்
  $\Subst{M''}{R''}{y}$ வரையறுக்கப்பட்டதாகவும் அமையும் மற்றொரு ஜோடி இருந்தால்,
  $\Subst{M'}{R'}{y} \aeq \Subst{M''}{R''}{y}$.
  % SOURCE CORRECTION TA-LAM-019: the second-pair premise must constrain R'' and its own substitution, not repeat the first pair's substitution.
\end{cor}

\begin{proof}
  \olref{thm:sub} இலிருந்து உடனடியாகக் கிடைக்கிறது.
\end{proof}

\end{document}
